% !PRINTED 670
\weberpage{670}
%<<CONT
in which $a, b, c, d$ are constants whose determinant $ad - bc$ is different from zero, then the branch polygons $\mathfrak{Z}_z$ and $\mathfrak{Z}_{z'}$ are identical, and hence $w_z = w_{z'}$ also; for if at a point $z-z_0$ or $1:z$ becomes zero in the $e$th order, then at the same point

\[
z' - z_0' = \dfrac{(ad-bc)(z-z_0)}{(cz+d)(cz_0+d)}
\]

or, when $z_0$ is infinite, so that $z_0' = a:c$:

\[
z' - z_0' = \dfrac{-(ad-bc)}{c(cz+d)},
\]

and when $z_0' = \infty$, so that $cz_0+d = 0$:

\[
\dfrac{1}{z'} = \dfrac{cz+d}{az+b}
\]

is likewise infinitely small in the $c$th order.
% [sic] printed $c$ten Ordnung, should be $e$ten: the order $e$ in which $z-z_0$ or $1:z$ becomes zero (stated above); $c$ is one of the constants $a, b, c, d$

If we apply this to $z' = 1:z$, we obtain the branch points at which $z$ has a finite value from the field discriminant $D_z$ according to § 184, (24), (28). There are then still lacking the branch points at which $z=\infty$, and hence $z'=0$, and these follow in the same way from the vanishing roots of the field discriminant $D_{z'}$. Accordingly we have the theorem:

5. \gesperrt{The branching number $w_z$ is equal to the degree of the field discriminant $D_z$, increased by the number of the vanishing roots of $D_{z'}$. ($z' = 1:z$).}

\begin{center}
\textbf{§ 191. Polygon quotients and polygon classes.}
\end{center}

A function $\eta$ in $\Omega$ has an order number different from zero only at a finite number of points. The sum of the positive order numbers is just as great as the sum of the negative ones, namely equal to the order of the function $\eta$ (§ 189, 4.). If the order numbers of $\eta$ are known for every point $\mathfrak{P}$, then the function $\eta$ is thereby determined up to a constant factor. For if a second function $\eta'$ has the same order numbers everywhere, then $\eta:\eta'$ has the order number $0$ everywhere and is therefore constant (§ 189, 4.).

We form a polygon $\mathfrak{A}$, into which we take up every point at which $\eta$ has a positive order number as often as this order number indicates, and a second polygon $\mathfrak{B}$, into which we
%CONT>>
